\chapter{Introducción}

En diversas situaciones, cuando se intenta llevar la teoría a la práctica, y sobre todo cuando esta se desarrolla en el mundo físico, es necesario considerar una gran cantidad de situaciones distintas.

Esto ocurre al estudiar una epidemia. Aunque existen muchos modelos matemáticos que describen su evolución y comportamiento, es necesario incorporar cada vez más variables para estimar de forma más precisa la situación real a la que se enfrenta la población.

La pandemia de COVID-19 es un claro ejemplo de ello. Si se tuviera una idea fiable de cómo se desarrolla una epidemia, se podrían diseñar políticas públicas realmente útiles.

Esta es una tarea difícil, ya que en el modelo SIRD (susceptibles, infectados, recuperados y fallecidos) normalmente se toman los parámetros de la epidemia como constantes, aun cuando estos pueden cambiar con el tiempo debido a factores como las medidas sanitarias, el comportamiento de la población o la aparición de nuevas variantes.

\section{Objetivos}

Estimar los parámetros $\beta$, $\gamma$ y $\delta$, (tasa de contagio, tasa de recuperación y tasa de mortalidad, respectivamente) dependientes del tiempo, del modelo SIRD mediante el método de ensamble de los filtros de Kalman, a partir de datos sobre COVID-19 en México.

\section{Justificación}

Los filtros de Kalman de ensamble (EnKF, por sus siglas en inglés) son una técnica de asimilación de datos diseñada para estimar estados y parámetros en sistemas dinámicos, especialmente cuando estos son no lineales y están sujetos a incertidumbre. A diferencia del filtro de Kalman clásico, el EnKF utiliza un conjunto de simulaciones (ensamble) para representar la distribución de probabilidad del sistema, lo que permite capturar de manera más flexible la variabilidad y el ruido en los datos. Este enfoque combina las predicciones del modelo con observaciones reales mediante actualizaciones secuenciales, ajustando continuamente las estimaciones conforme se dispone de nueva información.

